\chapter{Body Bias Trimming Strategy}
\label{chap:trimming}

The slicer-based quantisation stage introduced in Chapter~\ref{chap:async}
defines each decision threshold exclusively through the PMOS-to-NMOS sizing
ratio of a CMOS inverter, as expressed by equation~(\ref{eq:trip_point}). This
makes the converter linearity directly dependent on the effective drive
strengths of the two device types, and therefore directly exposed to
fabrication process variability. This chapter describes the offline trimming
strategy adopted to correct the resulting deterministic trip-point deviations
across process corners. The technique builds on the static, bulk-driven
calibration approaches reviewed in Section~\ref{sec:calibration}: it adjusts
the threshold voltages of the slicer transistors through the body effect, in
the same spirit as the bulk-driven reference-ladder trimming of
Yao~\cite{yao_bulk_2011} discussed in Section~\ref{subsec:ladder_trimming},
but applied here to the bulk terminals of the slicer bank and indexed by the
identified process corner rather than by a per-comparator offset measurement.
The scope of this correction is deterministic inter-corner variation only:
random intra-die mismatch, governed by Pelgrom's
law~\cite{pelgrom_1989}, is not addressed by the corner-indexed table and is
discussed as future work in Chapter~\ref{chap:Conclusions}.

\section{Choice of Body Biasing Among the Reviewed Techniques}

Among the trimming techniques surveyed in Section~\ref{sec:calibration},
several alternatives were considered before adopting body biasing. Internal
resistive loading, in which a digitally switched resistor bank modulates the
load of each comparator branch, is well suited to differential comparator
stages but does not map naturally onto the slicer inverters of the proposed
architecture, which have no explicit load element to modulate; introducing one
would add capacitance to the only high-speed node of each slicer and degrade
its switching delay. Dynamic offset-cancellation schemes such as auto-zeroing
were excluded for the reasons detailed in Section~\ref{subsec:static_trimming}:
they require periodic clock phases, which would reintroduce into the converter
precisely the clocked activity that the asynchronous architecture is designed
to eliminate.

Bulk-driven adjustment, by contrast, satisfies all the constraints of the
event-driven operation simultaneously. The bulk terminals of the slicer
transistors are not part of the signal path, so driving them from a static
bias source adds no capacitance to the switching nodes and consumes no dynamic
power. Once the correction voltages are set, they are held by purely DC
conduction paths, preserving the idle behaviour of the converter when the
input is inactive. Finally, because the NMOS and PMOS bulk potentials act on
the pull-down and pull-up networks independently, they provide exactly the two
degrees of freedom needed to counteract both the symmetric (SS, FF) and the
asymmetric (FS, SF) corner perturbations characterised in
Section~\ref{sec:corner_dependence}.

\section{Function of NMOS and PMOS Transistors Under Different Body Bias}

The threshold voltage of a MOSFET is not solely determined by the device
geometry and doping profile but is also a function of the potential applied
to its body terminal relative to the source. This dependence, commonly referred
to as the body effect or substrate bias effect, is described for an NMOS
transistor by

\begin{equation}
    V_{tn} = V_{tn0} + \gamma_n
    \left( \sqrt{2\phi_{F,n} + V_{SB}} - \sqrt{2\phi_{F,n}} \right),
    \label{eq:body_effect_nmos}
\end{equation}

\noindent
where $V_{tn0}$ is the zero-bias threshold voltage, $\gamma_n$ is the body
effect coefficient, $\phi_{F,n}$ is the Fermi potential, and $V_{SB}$ is the
source-to-body voltage. The analogous expression for a PMOS transistor,
accounting for the sign conventions appropriate to a p-channel device,
takes the form

\begin{equation}
    |V_{tp}| = |V_{tp0}| + \gamma_p
    \left( \sqrt{2\phi_{F,p} + |V_{SB}|} - \sqrt{2\phi_{F,p}} \right),
    \label{eq:body_effect_pmos}
\end{equation}

\noindent
where the relevant source-to-body voltage for the PMOS is determined by the
potential difference between the source and the N-well. In both cases, an
increase in $|V_{SB}|$ raises the magnitude of the threshold voltage, a
condition referred to as reverse body biasing (RBB). Conversely, reducing
$|V_{SB}|$ below its nominal value or, in technologies that permit it,
forward-biasing the body junction lowers $|V_t|$, a condition known as
forward body biasing (FBB).

A direct consequence for the slicer inverters in the proposed architecture is
that each cell's trip point, given by equation~(\ref{eq:trip_point}),
depends on the threshold voltages of both transistors composing the inverter.
By substituting the body-bias-dependent expressions for $V_{tn}$ and $|V_{tp}|$
into the inverter threshold equation and expanding to first order, it is
apparent that a positive shift in $V_{tn}$ achieved by applying RBB to the
NMOS weakens the pull-down network relative to the pull-up network and
therefore raises $V_M$, while an equivalent increase in $|V_{tp}|$ weakens
the pull-up network and lowers $V_M$. This asymmetry between the NMOS and
PMOS body bias responses provides two independent degrees of freedom for
adjusting the trip point of each slicer: the NMOS bulk voltage controls the
effective strength of the pull-down path, while the PMOS bulk voltage
(equivalently, the N-well potential) controls the pull-up path. Applying
corrections simultaneously to both body terminals therefore allows independent
tuning of the trip point magnitude and of the symmetry of the transition
characteristic, providing a richer correction space than either terminal alone
could offer.

\section{Process Corner Dependence and Motivation for Trimming}
\label{sec:corner_dependence}

In a deep-submicron CMOS process, fabrication variability causes the threshold
voltages and carrier mobilities of the NMOS and PMOS devices to deviate from
their nominal (typical-typical, TT) values. The standard process corners used
to bound this variability, namely slow-slow (SS), fast-fast (FF), fast-NMOS
slow-PMOS (FS) and slow-NMOS fast-PMOS (SF), capture the four extreme
combinations of NMOS and PMOS behaviour and represent the conditions under
which circuit performance is most likely to degrade.

For the slicer-based thermometer coder described in Chapter~\ref{chap:async},
the impact of process variation is particularly significant because the
thermometer code relies on the precise relative ordering of the inverter trip
points across the output bank. At the typical corner, the ratios $r_k$
assigned to each slicer place the trip points at the intended positions
relative to the reference level of $0.4$~V, and the resulting differential
non-linearity (DNL) and integral non-linearity (INL) of the converter remain
within specification. At non-typical corners, however, the effective
transistor strengths deviate from their nominal values in a manner that is not
uniform across the bank: since the PMOS and NMOS devices within each slicer
are affected differently depending on the corner, the trip point of each
inverter shifts by a cell-dependent amount, breaking the intended monotonic
spacing and introducing non-linearities that manifest as elevated DNL and INL.

At the SS corner, both NMOS and PMOS threshold voltages increase, reducing the
drive strength of both pull-down and pull-up networks. Although the symmetric
nature of this degradation might suggest a self-compensating effect on the
trip point, the body effect coefficients $\gamma_n$ and $\gamma_p$ and the
nominal threshold values are in general not matched between the two device
types, so the net shift in $V_M$ is non-zero and, more importantly,
cell-dependent due to the different sizing ratios across the slicer bank. The
FF corner produces the complementary effect: reduced threshold voltages
increase the drive strength of both device types, again shifting the trip
points by amounts that vary across the bank. The FS and SF corners are
qualitatively different in that they introduce an asymmetric perturbation.
At FS the NMOS is faster than nominal while the PMOS is slower, which biases
all slicers toward a lower trip point and compresses the effective decision
levels toward the bottom of the input range; at SF the complementary imbalance
expands the decision levels toward the top. In all four non-typical corners,
the loss of linearity was found to cause DNL excursions beyond $0.5$~LSB,
which makes the converter non-monotonic under these conditions.

\section{Offline Trimming Strategy via Body Bias Adjustment}

To restore the linearity of the thermometer coder across process corners
without modifying the core circuit topology, an offline trimming strategy
based on body bias adjustment was adopted. The principle of the approach is
to apply distinct bulk voltages to the NMOS and PMOS transistors of the output
slicer stage independently for each device type in a manner calibrated to
shift the trip points back toward their nominal positions. Since both the NMOS
bulk voltage $V_{BN}$ and the PMOS bulk voltage $V_{BP}$ contribute
independently to the trip point through equations (\ref{eq:body_effect_nmos})
and (\ref{eq:body_effect_pmos}), respectively, adjusting both simultaneously
corrects both the common-mode shift of the trip points, which affects the
DNL, and the differential imbalance between NMOS-dominated and PMOS-dominated
corners, which affects the INL. The trimming is offline in the sense that the
bulk voltages are determined once, from the identified process corner, and
held constant during normal operation; no continuous adaptive feedback is
involved.

\subsection{Determination of the Correction Pairs}

For each of the four non-typical corners, a specific pair of values
$(V_{BN},\, V_{BP})$ was determined such that the DNL of the converted
thermometer code falls below $0.5$~LSB across the full input range. The pairs
were obtained by sweeping the two bulk voltages over the range permitted by
the technology, extracting the trip points of all slicers at each bias
combination, and selecting the combination that minimises the worst-case
deviation of the measured trip points from their nominal TT positions. At the
SS corner, where both device types exhibit elevated threshold voltages and
therefore reduced drive strength, forward body biasing is applied to both the
NMOS and PMOS, reducing $V_{tn}$ and $|V_{tp}|$ toward their nominal values
and thereby restoring the intended trip point positions. At the FF corner,
reverse body biasing is applied to both types to counteract the excessive
drive strength resulting from the low threshold voltages. The FS and SF
corners require asymmetric corrections. At FS, a stronger reverse bias on the
NMOS bulk weakens the pull-down network to compensate for the fast-NMOS
tendency, while the PMOS bulk is forward-biased to recover part of the lost
pull-up strength; the SF corner is handled symmetrically, with the dominant
correction applied to the PMOS bulk. The resulting
correction table is summarised in Table~\ref{tab:trimming_vectors}.

\begin{table}[htbp]
    \centering
    \caption{Body bias correction vectors determined for each process corner.}
    \label{tab:trimming_vectors}
    \renewcommand{\arraystretch}{1.3}
    \begin{tabular}{lccc}
        \toprule
        \textbf{Corner} & \textbf{$V_{BN}$ (V)} & \textbf{$V_{BP}$ (V)} & \textbf{Bias type (NMOS / PMOS)} \\
        \midrule
        TT &  &  & nominal / nominal \\
        SS &  &  & FBB / FBB \\
        FF &  &  & RBB / RBB \\
        FS &  &  & RBB / FBB \\
        SF &  &  & FBB / RBB \\
        \bottomrule
    \end{tabular}
\end{table}

\subsection{Verilog-A Behavioural Implementation}

The trimming was implemented in simulation using a Verilog-A behavioural
model. The model monitors the process corner identifier and, based on the
detected corner, drives $V_{BN}$ and $V_{BP}$ to the pre-computed correction
values of Table~\ref{tab:trimming_vectors}. The two bias voltages are exposed
as electrical outputs of the behavioural block and routed to the bulk
terminals of the NMOS devices and to the N-well contacts of the PMOS devices
of the slicer bank, so that the correction acts on the transistor-level
netlist without requiring any modification to it. This approach allows the
trimming to be evaluated across all corners within the same simulation
framework as the nominal design. The Verilog-A abstraction also provides a
clean interface for extending the trimming table to intermediate corner
points or to temperature-corner combinations should a finer correction
granularity be required in future iterations. The effect of the body bias
corrections on the DNL and INL of the converter is presented in
Chapter~\ref{chap:Benchmarking}.
